\documentclass[aspectratio=169,11pt]{beamer}
\usetheme{default}
\usecolortheme{dove}
\usepackage{amsmath,amssymb,graphicx,tikz,pgfplots,booktabs}
\usetikzlibrary{arrows.meta,positioning,calc}
\pgfplotsset{compat=1.18}
\definecolor{DeckNavy}{HTML}{17324D}
\definecolor{DeckBlue}{HTML}{2D70B3}
\definecolor{DeckGreen}{HTML}{6F9940}
\definecolor{DeckRed}{HTML}{D95F4A}
\definecolor{DeckGold}{HTML}{D4A72C}
\definecolor{DeckGray}{HTML}{65727E}
\definecolor{DeckLight}{HTML}{F2F5F7}
\setbeamercolor{normal text}{fg=DeckNavy,bg=white}
\setbeamercolor{frametitle}{fg=DeckNavy,bg=white}
\setbeamercolor{title}{fg=DeckNavy}
\setbeamercolor{structure}{fg=DeckBlue}
\setbeamerfont{title}{size=\LARGE,series=\bfseries}
\setbeamerfont{frametitle}{size=\Large,series=\bfseries}
\setbeamertemplate{navigation symbols}{}
\setbeamertemplate{frametitle}{\vspace{0.25em}\hspace{0.15em}\insertframetitle\par
  \vspace{-0.35em}\hspace{0.15em}{\color{DeckGreen}\rule{1.3cm}{1.5pt}}}
\setbeamertemplate{footline}{\hfill\color{DeckGray}\scriptsize\insertframenumber\hspace{0.5cm}\vspace{0.2cm}}
\setbeamersize{text margin left=0.65cm,text margin right=0.65cm}
\hyphenpenalty=10000
\exhyphenpenalty=10000
\hypersetup{pdftitle={SGM Graph Filter Progress Update},pdfauthor={Harshil Shah},pdfsubject={MLP Implementation, Training, and Results}}
\newenvironment{canvas}{\begin{tikzpicture}[remember picture,overlay,shift={(current page.north west)},x=1cm,y=-1cm]}{\end{tikzpicture}}
\tikzset{txt/.style={anchor=north west,align=left,inner sep=0,text=DeckNavy},
  smalltxt/.style={txt,font=\small},
  labeltxt/.style={txt,font=\scriptsize,text=DeckGray},
  arr/.style={-{Latex[length=2mm]},thick,DeckGreen}}
\newcommand{\code}[1]{\texttt{\detokenize{#1}}}
\newcommand{\callout}[1]{\node[txt,fill=DeckNavy,text=white,rounded corners=2pt,
  inner sep=0.19cm,text width=13.8cm,align=center,font=\small\bfseries] at (0.9,7.20) {#1};}
\newcommand{\source}[1]{\node[labeltxt,text width=14.5cm,font=\tiny] at (0.75,8.12) {#1};}
\pgfplotsset{deckaxis/.style={axis lines=left,axis line style={DeckGray},
  tick label style={font=\scriptsize,text=DeckGray},label style={font=\small,text=DeckNavy},
  title style={font=\small\bfseries,text=DeckNavy},
  grid=major,grid style={DeckGray!14},tick style={DeckGray},
  legend style={draw=none,font=\scriptsize,fill=none},clip=true}}



\usepackage{listings}
\lstset{language=Python,basicstyle=\ttfamily\fontsize{9}{11}\selectfont,
 keywordstyle=\color{DeckBlue},commentstyle=\color{DeckGray},stringstyle=\color{DeckGreen},
 columns=fullflexible,keepspaces=true,showstringspaces=false,breaklines=true,
 aboveskip=0pt,belowskip=0pt}
\newcommand{\stepcontent}[4]{
\node[txt,font=\large\bfseries,text width=14cm] at (0.85,1.55) {#1};
\node[smalltxt,text width=6.15cm] at (0.85,2.45) {#2};
\node[txt,text width=7.25cm] at (7.75,3.2) {\begin{minipage}{7.25cm}\hyphenpenalty=0\exhyphenpenalty=0\lstinputlisting[linewidth=7.25cm]{snippets/#3.py}\end{minipage}};
\node[labeltxt,text width=7.25cm] at (7.75,2.35) {#4};}
\begin{document}
\begin{frame}[plain,noframenumbering]
\begin{canvas}
\node[txt,font=\LARGE\bfseries,text width=13.7cm] at (0.9,2.05) {SGM Graph Filter\\Progress Update};
\draw[DeckGreen,line width=1.5pt] (0.9,3.7) -- (2.4,3.7);
\node[txt,font=\large,text width=14cm] at (0.9,4.18) {MLP Implementation, Training, and Results};
\node[smalltxt] at (0.9,5.45) {Harshil Shah\\[0.18cm]Mission San Jose High School\\[0.18cm]October 19, 2026};
\end{canvas}
\end{frame}
\begin{frame}[t]{MLP Implementation}
\begin{canvas}
\only<1>{\stepcontent{1) Compute inputs offline}{For each subject and frequency, we compute the delayed graph and its eigenmodes once, then save them for training.\\[0.4cm]We use 36 subjects and 40 frequencies from 2--45 Hz. The graph has 86 nodes, with 68 cortical regions used for evaluation.}{graphs}{\code{training/data.py}\\[0.12cm]\code{generate_training_data}}}
\only<2>{\stepcontent{2) Build the MLP}{\textbf{Input:} normalized frequency, $f/45$.\\[0.35cm]\textbf{Output:} two real values, $a$ and $b$, which form the complex filter $a+bi$.\\[0.35cm]\textbf{Hyperparameters:} number of hidden layers, units per layer, and activation (Tanh or Sigmoid).\\[0.35cm]We vary these settings across experiments and learn only the MLP weights.}{layers}{\code{training/model.py}\\[0.12cm]\code{FeGraphMLP.__init__}}}
\only<3>{\stepcontent{3) Predict the complex filter}{We scale frequency before passing it into the MLP:\\[0.25cm]$x=\dfrac{\omega}{2\pi\cdot45}=\dfrac{f}{45}$.\\[0.35cm]The two outputs give the real and imaginary parts:\\[0.25cm]$F_\theta(\omega)=a_\theta(x)+i b_\theta(x)$.\\[0.35cm]Within each fold, we use the same filter for every subject. Frequency is the only input.}{input}{\code{training/model.py}\\[0.12cm]\code{FeGraphMLP.forward}}}
\only<4>{\stepcontent{4) Pass the filter into the SGM}{We pass $a+bi$ as \code{Fe_graph}, replacing $F_e$ only in graph propagation.\\[0.4cm]The existing SGM code handles the denominator safeguard, combines the graph modes, and selects the cortical regions.\\[0.4cm]This gives 68 regional spectra at 40 frequencies. The local response and fitted SGM parameters stay fixed.}{downstream}{\code{training/spectra.py}\\[0.12cm]\code{predict_spectra}}}
\end{canvas}
\note{Code excerpts follow the current repository. Line wrapping and explicitly marked omissions are for readability.}
\end{frame}

\begin{frame}[t]{Training and Evaluation}
\begin{canvas}
\only<1>{\stepcontent{1) Leave one subject out}{One \textbf{fold} is one train-and-test round.\\[0.3cm]\textbf{Fold 1:} train on subjects 2--36; test on subject 1.\\[0.3cm]\textbf{Fold 2:} start again, train on everyone except subject 2; test on subject 2.\\[0.3cm]Repeat until every subject has been held out once. Each round starts with fresh weights and a fresh optimizer, using seed 7.}{init}{\code{training/train.py}\\[0.12cm]\code{train_filter}}}
\only<2>{\stepcontent{2) Prepare spectra for comparison}{Within a fold, we compare each training subject\textquotesingle s predicted spectrum with their measured MEG spectrum.\\[0.3cm]Before calculating correlation, we process both in the same way:\\[0.2cm]1. Take the magnitude.\\[0.12cm]2. Smooth across neighboring frequencies.\\[0.12cm]3. Take the square root.\\[0.12cm]4. Subtract the mean across frequencies.}{transform}{\code{training/spectra.py}\\[0.12cm]\code{transform_spectra}}}
\only<3>{\stepcontent{3) Calculate the loss}{We calculate correlation across frequency for each region, then average the 68 regions:\\[0.3cm]$r_s=\dfrac1{68}\displaystyle\sum_{r=1}^{68}\mathrm{corr}_f(\widehat x_{s,r},y_{s,r})$.\\[0.35cm]The loss averages the 35 training subjects:\\[0.3cm]$\mathcal L=1-\dfrac1{35}\displaystyle\sum_{s\in\mathrm{train}}r_s$.\\[0.35cm]The held-out subject does not enter this loss.}{loss}{\code{training/train.py}\\[0.12cm]\code{train_filter}}}
\only<4>{\stepcontent{4) Update the MLP weights}{Each epoch makes one Adam update using all 35 training subjects. We record the loss before that update.\\[0.4cm]We run for a fixed 100, 500, or 5,000 epochs.\\[0.4cm]The initial learning rate is $10^{-3}$; the three follow-up runs use $3\times10^{-4}$.}{update}{\code{training/train.py}\\[0.12cm]\code{train_filter} (checks omitted)}}
\only<5>{\stepcontent{5) Evaluate the held-out subject}{After training, we apply the learned filter to the subject left out of that fold and compare the predicted spectra with their MEG.\\[0.35cm]We use the same spectrum processing and correlation calculation. This score never updates the MLP or selects a checkpoint.\\[0.35cm]After all 36 folds, each subject has one held-out score. We average those 36 scores to report performance.}{evaluate}{\code{training/train.py}\\[0.12cm]\code{train_filter}}}
\end{canvas}
\note{Code excerpts follow the current repository. Line wrapping and explicitly marked omissions are for readability.}
\end{frame}

\begin{frame}[t]{Experiments --- Initial comparison}\begin{canvas}\node[txt,text width=14cm,align=center] at (0.85,1.55) {\footnotesize\renewcommand{\arraystretch}{1.04}\begin{tabular}{r l l r r l}\toprule Run & Hidden layers & Activation & Epochs & Learning rate & Status\\\midrule0 & [8] & Tanh & 100 & 0.001 & Complete \\
1 & [8] & Sigmoid & 100 & 0.001 & Complete \\
2 & [8] & Tanh & 500 & 0.001 & Complete \\
3 & [8] & Sigmoid & 500 & 0.001 & Complete \\
4 & [32] & Tanh & 100 & 0.001 & Complete \\
5 & [32] & Sigmoid & 100 & 0.001 & Complete \\
6 & [32] & Tanh & 500 & 0.001 & Complete \\
7 & [32] & Sigmoid & 500 & 0.001 & Complete \\
8 & [32, 32] & Tanh & 100 & 0.001 & Complete \\
9 & [32, 32] & Sigmoid & 100 & 0.001 & Complete \\
10 & [32, 32] & Tanh & 500 & 0.001 & Complete \\
11 & [32, 32] & Sigmoid & 500 & 0.001 & Complete \\\bottomrule\end{tabular}};\node[labeltxt,text width=14cm,align=center] at (0.85,7.65) {};\end{canvas}\end{frame}

\begin{frame}[t]{Experiments --- Larger model and follow-ups}\begin{canvas}\node[txt,text width=14cm,align=center] at (0.85,1.55) {\footnotesize\renewcommand{\arraystretch}{1.04}\begin{tabular}{r l l r r l}\toprule Run & Hidden layers & Activation & Epochs & Learning rate & Status\\\midrule12 & $5\times128$ & Tanh & 5,000 & 0.001 & Complete \\
13 & $5\times128$ & Tanh & 5,000 & 0.0003 & Complete \\
14 & [32] & Tanh & 5,000 & 0.0003 & Complete \\
15 & [32, 32] & Tanh & 5,000 & 0.0003 & Complete \\\bottomrule\end{tabular}};\node[labeltxt,text width=14cm,align=center] at (0.85,7.65) {Runs 12 and 13 differ only in learning rate. Runs 13--15 use the same epoch budget.\\All three follow-ups are complete. Run 16 was outside this batch.};\end{canvas}\end{frame}

\begin{frame}[t]{Performance}\begin{canvas}\node[txt,text width=14cm,align=center] at (0.85,1.55) {\fontsize{9}{10}\selectfont\renewcommand{\arraystretch}{0.78}\begin{tabular}{r r r r}\toprule Run & Mean held-out $r$ & $\Delta r$ vs. analytical & Subjects improved\\\midrule Analytical & 0.8524 & --- & ---\\\midrule0 & 0.8163 & -0.0361 & 4/36\\
1 & 0.8170 & -0.0354 & 4/36\\
2 & 0.8511 & -0.0013 & 20/36\\
3 & 0.8355 & -0.0169 & 16/36\\
4 & 0.8377 & -0.0147 & 16/36\\
5 & 0.8280 & -0.0244 & 15/36\\
6 & 0.8562 & +0.0038 & 25/36\\
7 & 0.8303 & -0.0221 & 15/36\\
8 & 0.8371 & -0.0153 & 16/36\\
9 & 0.8276 & -0.0247 & 15/36\\
10 & 0.8511 & -0.0013 & 20/36\\
11 & 0.8377 & -0.0147 & 16/36\\
12 & 0.8605 & +0.0081 & 27/36\\
13 & 0.8620 & +0.0096 & 31/36\\
14 & 0.8577 & +0.0053 & 26/36\\
15 & 0.8599 & +0.0075 & 29/36\\\bottomrule\end{tabular}};\node[labeltxt,text width=14cm,align=center] at (0.85,8.1) {Exploratory single-seed comparisons; selecting a winner from these scores is model selection.};\end{canvas}\end{frame}

\begin{frame}[t]{Run 0 --- [8] Tanh}\begin{canvas}\node[smalltxt] at (0.85,1.55) {100 epochs; Adam 0.001; seed 7};\node[txt] at (0.35,2.2) {\includegraphics[width=7.45cm,height=4.45cm,keepaspectratio]{../figures/run_0_training_curves.png}};\node[txt] at (8.0,2.2) {\includegraphics[width=7.45cm,height=4.45cm,keepaspectratio]{../figures/run_0_spectra_across_subjects.png}};\node[smalltxt,text width=14cm] at (0.85,6.75) {Held-out $r$: \textbf{0.8163}\qquad $\Delta r$: -0.0361\qquad Improved: 4/36 subjects};\node[labeltxt,text width=14cm] at (0.85,7.45) {Left: training loss for each fold and their mean. Right: normalized spectral shape, mean $\pm$ SD across subjects.};\end{canvas}\note{Source: results/experiments/basic\_mlp\_0. The spectral plot averages normalized regions within subject, then subjects equally. SD is across subjects, not a confidence interval. Training loss is before each optimizer step.}\end{frame}

\begin{frame}[t]{Run 1 --- [8] Sigmoid}\begin{canvas}\node[smalltxt] at (0.85,1.55) {100 epochs; Adam 0.001; seed 7};\node[txt] at (0.35,2.2) {\includegraphics[width=7.45cm,height=4.45cm,keepaspectratio]{../figures/run_1_training_curves.png}};\node[txt] at (8.0,2.2) {\includegraphics[width=7.45cm,height=4.45cm,keepaspectratio]{../figures/run_1_spectra_across_subjects.png}};\node[smalltxt,text width=14cm] at (0.85,6.75) {Held-out $r$: \textbf{0.8170}\qquad $\Delta r$: -0.0354\qquad Improved: 4/36 subjects};\node[labeltxt,text width=14cm] at (0.85,7.45) {Left: training loss for each fold and their mean. Right: normalized spectral shape, mean $\pm$ SD across subjects.};\end{canvas}\note{Source: results/experiments/basic\_mlp\_1. The spectral plot averages normalized regions within subject, then subjects equally. SD is across subjects, not a confidence interval. Training loss is before each optimizer step.}\end{frame}

\begin{frame}[t]{Run 2 --- [8] Tanh}\begin{canvas}\node[smalltxt] at (0.85,1.55) {500 epochs; Adam 0.001; seed 7};\node[txt] at (0.35,2.2) {\includegraphics[width=7.45cm,height=4.45cm,keepaspectratio]{../figures/run_2_training_curves.png}};\node[txt] at (8.0,2.2) {\includegraphics[width=7.45cm,height=4.45cm,keepaspectratio]{../figures/run_2_spectra_across_subjects.png}};\node[smalltxt,text width=14cm] at (0.85,6.75) {Held-out $r$: \textbf{0.8511}\qquad $\Delta r$: -0.0013\qquad Improved: 20/36 subjects};\node[labeltxt,text width=14cm] at (0.85,7.45) {Left: training loss for each fold and their mean. Right: normalized spectral shape, mean $\pm$ SD across subjects.};\end{canvas}\note{Source: results/experiments/basic\_mlp\_2. The spectral plot averages normalized regions within subject, then subjects equally. SD is across subjects, not a confidence interval. Training loss is before each optimizer step.}\end{frame}

\begin{frame}[t]{Run 3 --- [8] Sigmoid}\begin{canvas}\node[smalltxt] at (0.85,1.55) {500 epochs; Adam 0.001; seed 7};\node[txt] at (0.35,2.2) {\includegraphics[width=7.45cm,height=4.45cm,keepaspectratio]{../figures/run_3_training_curves.png}};\node[txt] at (8.0,2.2) {\includegraphics[width=7.45cm,height=4.45cm,keepaspectratio]{../figures/run_3_spectra_across_subjects.png}};\node[smalltxt,text width=14cm] at (0.85,6.75) {Held-out $r$: \textbf{0.8355}\qquad $\Delta r$: -0.0169\qquad Improved: 16/36 subjects};\node[labeltxt,text width=14cm] at (0.85,7.45) {Left: training loss for each fold and their mean. Right: normalized spectral shape, mean $\pm$ SD across subjects.};\end{canvas}\note{Source: results/experiments/basic\_mlp\_3. The spectral plot averages normalized regions within subject, then subjects equally. SD is across subjects, not a confidence interval. Training loss is before each optimizer step.}\end{frame}

\begin{frame}[t]{Run 4 --- [32] Tanh}\begin{canvas}\node[smalltxt] at (0.85,1.55) {100 epochs; Adam 0.001; seed 7};\node[txt] at (0.35,2.2) {\includegraphics[width=7.45cm,height=4.45cm,keepaspectratio]{../figures/run_4_training_curves.png}};\node[txt] at (8.0,2.2) {\includegraphics[width=7.45cm,height=4.45cm,keepaspectratio]{../figures/run_4_spectra_across_subjects.png}};\node[smalltxt,text width=14cm] at (0.85,6.75) {Held-out $r$: \textbf{0.8377}\qquad $\Delta r$: -0.0147\qquad Improved: 16/36 subjects};\node[labeltxt,text width=14cm] at (0.85,7.45) {Left: training loss for each fold and their mean. Right: normalized spectral shape, mean $\pm$ SD across subjects.};\end{canvas}\note{Source: results/experiments/basic\_mlp\_4. The spectral plot averages normalized regions within subject, then subjects equally. SD is across subjects, not a confidence interval. Training loss is before each optimizer step.}\end{frame}

\begin{frame}[t]{Run 5 --- [32] Sigmoid}\begin{canvas}\node[smalltxt] at (0.85,1.55) {100 epochs; Adam 0.001; seed 7};\node[txt] at (0.35,2.2) {\includegraphics[width=7.45cm,height=4.45cm,keepaspectratio]{../figures/run_5_training_curves.png}};\node[txt] at (8.0,2.2) {\includegraphics[width=7.45cm,height=4.45cm,keepaspectratio]{../figures/run_5_spectra_across_subjects.png}};\node[smalltxt,text width=14cm] at (0.85,6.75) {Held-out $r$: \textbf{0.8280}\qquad $\Delta r$: -0.0244\qquad Improved: 15/36 subjects};\node[labeltxt,text width=14cm] at (0.85,7.45) {Left: training loss for each fold and their mean. Right: normalized spectral shape, mean $\pm$ SD across subjects.};\end{canvas}\note{Source: results/experiments/basic\_mlp\_5. The spectral plot averages normalized regions within subject, then subjects equally. SD is across subjects, not a confidence interval. Training loss is before each optimizer step.}\end{frame}

\begin{frame}[t]{Run 6 --- [32] Tanh}\begin{canvas}\node[smalltxt] at (0.85,1.55) {500 epochs; Adam 0.001; seed 7};\node[txt] at (0.35,2.2) {\includegraphics[width=7.45cm,height=4.45cm,keepaspectratio]{../figures/run_6_training_curves.png}};\node[txt] at (8.0,2.2) {\includegraphics[width=7.45cm,height=4.45cm,keepaspectratio]{../figures/run_6_spectra_across_subjects.png}};\node[smalltxt,text width=14cm] at (0.85,6.75) {Held-out $r$: \textbf{0.8562}\qquad $\Delta r$: +0.0038\qquad Improved: 25/36 subjects};\node[labeltxt,text width=14cm] at (0.85,7.45) {Left: training loss for each fold and their mean. Right: normalized spectral shape, mean $\pm$ SD across subjects.};\end{canvas}\note{Source: results/experiments/basic\_mlp\_6. The spectral plot averages normalized regions within subject, then subjects equally. SD is across subjects, not a confidence interval. Training loss is before each optimizer step.}\end{frame}

\begin{frame}[t]{Run 7 --- [32] Sigmoid}\begin{canvas}\node[smalltxt] at (0.85,1.55) {500 epochs; Adam 0.001; seed 7};\node[txt] at (0.35,2.2) {\includegraphics[width=7.45cm,height=4.45cm,keepaspectratio]{../figures/run_7_training_curves.png}};\node[txt] at (8.0,2.2) {\includegraphics[width=7.45cm,height=4.45cm,keepaspectratio]{../figures/run_7_spectra_across_subjects.png}};\node[smalltxt,text width=14cm] at (0.85,6.75) {Held-out $r$: \textbf{0.8303}\qquad $\Delta r$: -0.0221\qquad Improved: 15/36 subjects};\node[labeltxt,text width=14cm] at (0.85,7.45) {Left: training loss for each fold and their mean. Right: normalized spectral shape, mean $\pm$ SD across subjects.};\end{canvas}\note{Source: results/experiments/basic\_mlp\_7. The spectral plot averages normalized regions within subject, then subjects equally. SD is across subjects, not a confidence interval. Training loss is before each optimizer step.}\end{frame}

\begin{frame}[t]{Run 8 --- [32, 32] Tanh}\begin{canvas}\node[smalltxt] at (0.85,1.55) {100 epochs; Adam 0.001; seed 7};\node[txt] at (0.35,2.2) {\includegraphics[width=7.45cm,height=4.45cm,keepaspectratio]{../figures/run_8_training_curves.png}};\node[txt] at (8.0,2.2) {\includegraphics[width=7.45cm,height=4.45cm,keepaspectratio]{../figures/run_8_spectra_across_subjects.png}};\node[smalltxt,text width=14cm] at (0.85,6.75) {Held-out $r$: \textbf{0.8371}\qquad $\Delta r$: -0.0153\qquad Improved: 16/36 subjects};\node[labeltxt,text width=14cm] at (0.85,7.45) {Left: training loss for each fold and their mean. Right: normalized spectral shape, mean $\pm$ SD across subjects.};\end{canvas}\note{Source: results/experiments/basic\_mlp\_8. The spectral plot averages normalized regions within subject, then subjects equally. SD is across subjects, not a confidence interval. Training loss is before each optimizer step.}\end{frame}

\begin{frame}[t]{Run 9 --- [32, 32] Sigmoid}\begin{canvas}\node[smalltxt] at (0.85,1.55) {100 epochs; Adam 0.001; seed 7};\node[txt] at (0.35,2.2) {\includegraphics[width=7.45cm,height=4.45cm,keepaspectratio]{../figures/run_9_training_curves.png}};\node[txt] at (8.0,2.2) {\includegraphics[width=7.45cm,height=4.45cm,keepaspectratio]{../figures/run_9_spectra_across_subjects.png}};\node[smalltxt,text width=14cm] at (0.85,6.75) {Held-out $r$: \textbf{0.8276}\qquad $\Delta r$: -0.0247\qquad Improved: 15/36 subjects};\node[labeltxt,text width=14cm] at (0.85,7.45) {Left: training loss for each fold and their mean. Right: normalized spectral shape, mean $\pm$ SD across subjects.};\end{canvas}\note{Source: results/experiments/basic\_mlp\_9. The spectral plot averages normalized regions within subject, then subjects equally. SD is across subjects, not a confidence interval. Training loss is before each optimizer step.}\end{frame}

\begin{frame}[t]{Run 10 --- [32, 32] Tanh}\begin{canvas}\node[smalltxt] at (0.85,1.55) {500 epochs; Adam 0.001; seed 7};\node[txt] at (0.35,2.2) {\includegraphics[width=7.45cm,height=4.45cm,keepaspectratio]{../figures/run_10_training_curves.png}};\node[txt] at (8.0,2.2) {\includegraphics[width=7.45cm,height=4.45cm,keepaspectratio]{../figures/run_10_spectra_across_subjects.png}};\node[smalltxt,text width=14cm] at (0.85,6.75) {Held-out $r$: \textbf{0.8511}\qquad $\Delta r$: -0.0013\qquad Improved: 20/36 subjects};\node[labeltxt,text width=14cm] at (0.85,7.45) {Left: training loss for each fold and their mean. Right: normalized spectral shape, mean $\pm$ SD across subjects.};\end{canvas}\note{Source: results/experiments/basic\_mlp\_10. The spectral plot averages normalized regions within subject, then subjects equally. SD is across subjects, not a confidence interval. Training loss is before each optimizer step.}\end{frame}

\begin{frame}[t]{Run 11 --- [32, 32] Sigmoid}\begin{canvas}\node[smalltxt] at (0.85,1.55) {500 epochs; Adam 0.001; seed 7};\node[txt] at (0.35,2.2) {\includegraphics[width=7.45cm,height=4.45cm,keepaspectratio]{../figures/run_11_training_curves.png}};\node[txt] at (8.0,2.2) {\includegraphics[width=7.45cm,height=4.45cm,keepaspectratio]{../figures/run_11_spectra_across_subjects.png}};\node[smalltxt,text width=14cm] at (0.85,6.75) {Held-out $r$: \textbf{0.8377}\qquad $\Delta r$: -0.0147\qquad Improved: 16/36 subjects};\node[labeltxt,text width=14cm] at (0.85,7.45) {Left: training loss for each fold and their mean. Right: normalized spectral shape, mean $\pm$ SD across subjects.};\end{canvas}\note{Source: results/experiments/basic\_mlp\_11. The spectral plot averages normalized regions within subject, then subjects equally. SD is across subjects, not a confidence interval. Training loss is before each optimizer step.}\end{frame}

\begin{frame}[t]{Run 12 --- $5\times128$ Tanh}\begin{canvas}\node[smalltxt] at (0.85,1.55) {5,000 epochs; Adam 0.001; seed 7};\node[txt] at (0.35,2.2) {\includegraphics[width=7.45cm,height=4.45cm,keepaspectratio]{../figures/run_12_training_curves.png}};\node[txt] at (8.0,2.2) {\includegraphics[width=7.45cm,height=4.45cm,keepaspectratio]{../figures/run_12_spectra_across_subjects.png}};\node[smalltxt,text width=14cm] at (0.85,6.75) {Held-out $r$: \textbf{0.8605}\qquad $\Delta r$: +0.0081\qquad Improved: 27/36 subjects};\node[labeltxt,text width=14cm] at (0.85,7.45) {Left: training loss for each fold and their mean. Right: normalized spectral shape, mean $\pm$ SD across subjects.};\end{canvas}\note{Source: results/experiments/basic\_mlp\_12. The spectral plot averages normalized regions within subject, then subjects equally. SD is across subjects, not a confidence interval. Training loss is before each optimizer step.}\end{frame}

\begin{frame}[t]{Run 13 --- $5\times128$ Tanh}\begin{canvas}\node[smalltxt] at (0.85,1.55) {5,000 epochs; Adam 0.0003; seed 7};\node[txt] at (0.35,2.2) {\includegraphics[width=7.45cm,height=4.45cm,keepaspectratio]{../figures/run_13_training_curves.png}};\node[txt] at (8.0,2.2) {\includegraphics[width=7.45cm,height=4.45cm,keepaspectratio]{../figures/run_13_spectra_across_subjects.png}};\node[smalltxt,text width=14cm] at (0.85,6.75) {Held-out $r$: \textbf{0.8620}\qquad $\Delta r$: +0.0096\qquad Improved: 31/36 subjects};\node[labeltxt,text width=14cm] at (0.85,7.45) {Left: training loss for each fold and their mean. Right: normalized spectral shape, mean $\pm$ SD across subjects.};\end{canvas}\note{Source: results/experiments/basic\_mlp\_13. The spectral plot averages normalized regions within subject, then subjects equally. SD is across subjects, not a confidence interval. Training loss is before each optimizer step.}\end{frame}

\begin{frame}[t]{Run 14 --- [32] Tanh}\begin{canvas}\node[smalltxt] at (0.85,1.55) {5,000 epochs; Adam 0.0003; seed 7};\node[txt] at (0.35,2.2) {\includegraphics[width=7.45cm,height=4.45cm,keepaspectratio]{../figures/run_14_training_curves.png}};\node[txt] at (8.0,2.2) {\includegraphics[width=7.45cm,height=4.45cm,keepaspectratio]{../figures/run_14_spectra_across_subjects.png}};\node[smalltxt,text width=14cm] at (0.85,6.75) {Held-out $r$: \textbf{0.8577}\qquad $\Delta r$: +0.0053\qquad Improved: 26/36 subjects};\node[labeltxt,text width=14cm] at (0.85,7.45) {Left: training loss for each fold and their mean. Right: normalized spectral shape, mean $\pm$ SD across subjects.};\end{canvas}\note{Source: results/experiments/basic\_mlp\_14. The spectral plot averages normalized regions within subject, then subjects equally. SD is across subjects, not a confidence interval. Training loss is before each optimizer step.}\end{frame}

\begin{frame}[t]{Run 15 --- [32, 32] Tanh}\begin{canvas}\node[smalltxt] at (0.85,1.55) {5,000 epochs; Adam 0.0003; seed 7};\node[txt] at (0.35,2.2) {\includegraphics[width=7.45cm,height=4.45cm,keepaspectratio]{../figures/run_15_training_curves.png}};\node[txt] at (8.0,2.2) {\includegraphics[width=7.45cm,height=4.45cm,keepaspectratio]{../figures/run_15_spectra_across_subjects.png}};\node[smalltxt,text width=14cm] at (0.85,6.75) {Held-out $r$: \textbf{0.8599}\qquad $\Delta r$: +0.0075\qquad Improved: 29/36 subjects};\node[labeltxt,text width=14cm] at (0.85,7.45) {Left: training loss for each fold and their mean. Right: normalized spectral shape, mean $\pm$ SD across subjects.};\end{canvas}\note{Source: results/experiments/basic\_mlp\_15. The spectral plot averages normalized regions within subject, then subjects equally. SD is across subjects, not a confidence interval. Training loss is before each optimizer step.}\end{frame}
\end{document}
